% AA-09 | Algebra abstrakcyjna | Iteracja V
\setlength{\parskip}{0pt}
\chapter{Pierścienie, ideały i homomorfizmy}
\label{aa:ch:pierscienie}

\section{Dlaczego jedno działanie już nie wystarcza?}
Grupa pozwala mówić o jednym działaniu i jego odwracalności.
Jednak w arytmetyce liczb całkowitych i w rachunku wielomianów
mamy \emph{dwa} działania: dodawanie i mnożenie. Zachowują się one
inaczej: każdą liczbę całkowitą można odjąć, lecz przez niezerową
liczbę całkowitą nie zawsze można podzielić z wynikiem całkowitym.
Z drugiej strony macierze można dodawać i mnożyć, ale mnożenie
macierzy na ogół nie jest przemienne. Szukamy struktury, która
obejmie wszystkie te przykłady i zachowa reguły rozdzielności.

\begin{przyklad}[Jedno równanie w trzech pierścieniach]
\label{aa:prz:pierscienie-rownanie}
Rozważmy równanie $x^2=1$. W pierścieniu $\mathbb Z$
ma ono dokładnie dwa rozwiązania: $1$ i $-1$. Istotnie,
$(x-1)(x+1)=0$; ponieważ z iloczynu dwóch liczb całkowitych
równego zeru wynika zerowość jednego czynnika, mamy $x=1$ albo
$x=-1$. W pierścieniu $\mathbb Z/8\mathbb Z$ rozwiązaniami
są klasy $[1],[3],[5],[7]$. Dla nieparzystej liczby $x=2k+1$
liczba $x^2-1=4k(k+1)$ jest podzielna przez $8$, gdyż jedna
z liczb $k,k+1$ jest parzysta. Klasy parzyste nie rozwiązują
równania, ponieważ ich kwadraty są parzyste. Wreszcie w pierścieniu
$\mathbb R[\varepsilon]/(\varepsilon^2)$ klasa wielomianu
$\varepsilon$ nie jest zerem (bo $\varepsilon\notin(\varepsilon^2)$,
co wynika ze stopni wielomianów), lecz jej kwadrat jest zerem.
Pojęcie pierścienia
musi zatem dopuszczać dzielniki zera i elementy nilpotentne.
\end{przyklad}

\section{Aksjomaty pierścienia i podstawowe konsekwencje}
\begin{definicja}[Pierścień z jedynką]
\label{aa:def:pierscien}
\emph{Pierścieniem z jedynką} nazywamy zbiór $R$ wyposażony
w działania $+:R\times R\to R$ oraz $\cdot:R\times R\to R$
takie, że:
\begin{enumerate}[label=\textup{(P\arabic*)}]
\item $(R,+)$ jest grupą abelową z elementem neutralnym $0_R$;
\item mnożenie jest łączne: $(ab)c=a(bc)$ dla wszystkich $a,b,c\in R$;
\item istnieje $1_R\in R$ takie, że $1_Ra=a1_R=a$ dla każdego $a\in R$;
\item $a(b+c)=ab+ac$ i $(a+b)c=ac+bc$ dla wszystkich $a,b,c\in R$.
\end{enumerate}
Pierścień jest \emph{przemienny}, jeśli $ab=ba$ dla wszystkich
$a,b\in R$. Nie zakładamy tego, jeśli nie zostało to wskazane.
Dopuszczamy pierścień zerowy $R=\{0\}$, w którym $1_R=0_R$;
we wszystkich twierdzeniach o dziedzinach i ciałach wykluczymy go jawnie.
\end{definicja}

\begin{lemat}[Rachunek znaków]
\label{aa:lem:pierscien-znaki}
Niech $R$ będzie pierścieniem z jedynką i niech $a,b\in R$.
Wtedy $a0_R=0_Ra=0_R$, $(-a)b=-(ab)=a(-b)$ oraz
$(-a)(-b)=ab$. Jeżeli $1_R=0_R$, to $R=\{0_R\}$.
\end{lemat}
\begin{proof}
Z rozdzielności otrzymujemy $a0_R=a(0_R+0_R)=a0_R+a0_R$.
Po dodaniu elementu odwrotnego do $a0_R$ dostajemy $a0_R=0_R$.
Analogicznie $0_Ra=0_R$. Następnie
$(-a)b+ab=((-a)+a)b=0_Rb=0_R$, więc $(-a)b=-(ab)$.
Z kolei $a(-b)+ab=a((-b)+b)=0_R$, zatem $a(-b)=-(ab)$.
Stosując pierwszą z tych równości do $-b$, dostajemy
$(-a)(-b)=-(a(-b))=-(-(ab))=ab$.
Jeżeli $1_R=0_R$, to każdy $a\in R$ spełnia
$a=1_Ra=0_Ra=0_R$, co dowodzi ostatniego zdania.
\end{proof}


\section{Konstrukcja pierścienia liczb całkowitych z liczb naturalnych}
\label{aa:sec:konstrukcja-z}
Wcześniej korzystaliśmy z liczb całkowitych jako znanego przykładu
pierścienia. Możemy teraz uzasadnić ich istnienie, nie zakładając
z góry, że odejmowanie liczb naturalnych jest zawsze możliwe.
W tej sekcji $\mathbb N_0=\{0,1,2,\ldots\}$ jest półpierścieniem
zwykłych liczb naturalnych, to znaczy ma łączne i przemienne dodawanie
z zerem, łączne i przemienne mnożenie z jedynką oraz rozdzielność.
Korzystamy z własności skracania dodawania w $\mathbb N_0$:
$a+t=b+t$ pociąga $a=b$; wynika ona z definicji dodawania
liczb naturalnych i można ją wykazać indukcyjnie względem $t$.

\begin{definicja}[Formalne różnice liczb naturalnych]
\label{aa:def:formal-differences}
Na $\mathbb N_0\times\mathbb N_0$ określamy relację
\[
 (a,b)\sim(c,d)\quad\Longleftrightarrow\quad a+d=b+c.
\]
Zbiór klas tej relacji oznaczymy przez $Z$; klasę pary $(a,b)$
zapisujemy $[a,b]$. Para $(a,b)$ ma wyrażać \emph{formalną różnicę}
$a-b$, lecz zapis ten jest na razie tylko intuicją: liczby ujemne
nie zostały jeszcze zdefiniowane.
\end{definicja}

\begin{lemat}[Relacja formalnych różnic]
\label{aa:lem:formal-differences-equivalence}
Niech $\sim$ oznacza warunek $(a,b)\sim(c,d)$ wtedy i tylko wtedy,
  gdy $a+d=b+c$ dla par z $\mathbb N_0\times\mathbb N_0$.
  Jest to relacja równoważności.
\end{lemat}
\begin{proof}
Dla każdej pary $(a,b)$ zachodzi $a+b=b+a$, więc relacja
jest zwrotna. Jeśli $a+d=b+c$, to $c+b=d+a$,
więc jest symetryczna. Wreszcie z $(a,b)\sim(c,d)$
oraz $(c,d)\sim(e,f)$ otrzymujemy
$a+d=b+c$ i $c+f=d+e$. Dodając $f$ do pierwszej
równości, a następnie podstawiając drugą, mamy
$a+d+f=b+c+f=b+d+e$. Dzięki skracaniu składnika $d$
dostajemy $a+f=b+e$, czyli $(a,b)\sim(e,f)$.
\end{proof}

\begin{definicja}[Działania na formalnych różnicach]
\label{aa:def:z-operations}
Niech $Z=(\mathbb N_0\times\mathbb N_0)/\sim$, gdzie
$(a,b)\sim(c,d)$ wtedy i tylko wtedy, gdy $a+d=b+c$.
Dla wszystkich $a,b,c,d\in\mathbb N_0$ określamy na $Z$
działania wzorami
\[
 [a,b]+[c,d]=[a+c,b+d],\qquad
 [a,b]\,[c,d]=[ac+bd,ad+bc].
\]
Pierwszy odpowiada dodawaniu różnic, a drugi tożsamości
$(a-b)(c-d)=ac+bd-(ad+bc)$, użytej jedynie jako
wskazówka przy wyborze wzoru.
\end{definicja}

\begin{lemat}[Działania są dobrze określone]
\label{aa:lem:z-well-defined}
Niech $Z=(\mathbb N_0\times\mathbb N_0)/\sim$, gdzie
$(a,b)\sim(c,d)$ wtedy i tylko wtedy, gdy $a+d=b+c$.
Dodawanie i mnożenie klas zadane wzorami z
Definicji~\ref{aa:def:z-operations} nie zależą od wyboru reprezentantów.
\end{lemat}
\begin{proof}
Niech $(a,b)\sim(a',b')$, a więc $a+b'=b+a'$.
Dla dowolnej pary $(c,d)$ otrzymujemy
$(a+c)+(b'+d)=(b+d)+(a'+c)$, co dowodzi
niezależności dodawania od reprezentanta pierwszego składnika.
Analogicznie postępujemy z drugim składnikiem.

Dla mnożenia należy wykazać
$(ac+bd,ad+bc)\sim(a'c+b'd,a'd+b'c)$.
Warunek ten oznacza równość
\[
 (ac+bd)+(a'd+b'c)=(ad+bc)+(a'c+b'd).
\]
Lewą stronę przestawiamy na $(a+b')c+(b+a')d$,
a prawą na $(b+a')c+(a+b')d$; są one równe, bo
$a+b'=b+a'$. Dowodzi to niezależności od
reprezentanta pierwszego czynnika. Ponieważ podany wzór
na mnożenie jest symetryczny względem par czynników,
ten sam argument stosuje się do drugiego czynnika.
Po dwóch kolejnych zmianach reprezentantów otrzymujemy
niezależność obu działań jednocześnie.
\end{proof}

\begin{twierdzenie}[Pierścień liczb całkowitych]
\label{aa:tw:integers-as-ring}
Niech $Z=(\mathbb N_0\times\mathbb N_0)/\sim$, gdzie
$(a,b)\sim(c,d)$ wtedy i tylko wtedy, gdy $a+d=b+c$.
Niech dla wszystkich $a,b,c,d\in\mathbb N_0$ zachodzą wzory
$[a,b]+[c,d]=[a+c,b+d]$ oraz
$[a,b][c,d]=[ac+bd,ad+bc]$.
Wówczas $Z$ jest niezerowym pierścieniem przemiennym z jedynką.
Jego elementem zerowym
jest $[0,0]$, jedynką $[1,0]$, a przeciwnym do $[a,b]$
jest $[b,a]$. Odwzorowanie
\[
 \iota:\mathbb N_0\longrightarrow Z,\qquad \iota(n)=[n,0]
\]
jest różnowartościowe oraz zachowuje dodawanie,
mnożenie, zero i jedynkę. Pierścień $Z$ utożsamiamy
odtąd z pierścieniem $\mathbb Z$.
\end{twierdzenie}
\begin{proof}
Najpierw sprawdzamy strukturę addytywną. Łączność,
przemienność i neutralność $[0,0]$ wynikają z
odpowiednich reguł dodawania w $\mathbb N_0$.
Przy tym
$[a,b]+[b,a]=[a+b,a+b]=[0,0]$,
bo $(a+b)+0=(a+b)+0$. Każda klasa ma więc element
przeciwny i $(Z,+)$ jest grupą abelową.

Przemienność mnożenia wynika ze wzoru i przemienności
mnożenia w $\mathbb N_0$. Dla łączności porównajmy
współrzędne obu iloczynów trzech par. Pierwsza współrzędna
$(([a,b][c,d])[e,f])$ wynosi
$ace+bde+adf+bcf$, a druga
$acf+bdf+ade+bce$. W iloczynie
$[a,b]([c,d][e,f])$ otrzymujemy dokładnie te same
składniki, jedynie w innej kolejności. Oba wyniki są więc
równe już jako pary. Ponadto
$[a,b][1,0]=[a,b]$. Wreszcie rozdzielność
mnożenia względem dodawania wynika z rachunku
\[
 [a,b]([c,d]+[e,f])
 =[a(c+e)+b(d+f),\,a(d+f)+b(c+e)]
 =[a,b][c,d]+[a,b][e,f].
\]
Rozdzielność z prawej wynika z przemienności mnożenia.
Zatem mamy przemienny pierścień z jedynką.

Elementy $[1,0]$ i $[0,0]$ nie są równe, bo równość
ich klas wymagałaby $1+0=0+0$ w $\mathbb N_0$.
Odwzorowanie $\iota$ zachowuje oba działania i elementy
wyróżnione bezpośrednio ze wzorów. Jeśli
$[m,0]=[n,0]$, to $m+0=0+n$, zatem $m=n$;
$\iota$ jest różnowartościowe.
\end{proof}

\begin{twierdzenie}[Własność uniwersalna dopełnienia grupowego]
\label{aa:tw:z-grothendieck-universal}
Niech $\mathbb Z=(\mathbb N_0\times\mathbb N_0)/\sim$, gdzie
$(a,b)\sim(c,d)$ wtedy i tylko wtedy, gdy $a+d=b+c$,
oraz niech $\iota:\mathbb N_0\to\mathbb Z$ będzie dane wzorem
$\iota(n)=[n,0]$. Niech $A$ będzie dowolną grupą abelową
zapisywaną addytywnie,
a $f:\mathbb N_0\to A$ odwzorowaniem spełniającym
$f(m+n)=f(m)+f(n)$ i $f(0)=0$. Istnieje dokładnie jeden
homomorfizm grup abelowych $\overline f:\mathbb Z\to A$
taki, że $\overline f\circ\iota=f$. Jest nim
\[
 \overline f([a,b])=f(a)-f(b).
\]
Co więcej, dla każdego przemiennego pierścienia $R$
z jedynką jedyny homomorfizm pierścieni
$\Phi:\mathbb Z\to R$ ma postać
$\Phi([a,b])=a1_R-b1_R$.
\end{twierdzenie}
\begin{proof}
Jeśli $(a,b)\sim(c,d)$, to $a+d=b+c$.
Po zastosowaniu addytywności $f$ dostajemy
$f(a)+f(d)=f(b)+f(c)$, stąd
$f(a)-f(b)=f(c)-f(d)$. Wzór nie zależy więc
od reprezentanta. Dodawanie klas przechodzi na dodawanie
w $A$, zatem $\overline f$ jest homomorfizmem,
a $\overline f([n,0])=f(n)$. Każda klasa jest różnicą
$[a,0]-[b,0]$, dlatego każdy homomorfizm spełniający
ten warunek musi przyjmować podane wartości. To dowodzi
jednoznaczności.

Jeśli $R$ jest pierścieniem z jedynką, homomorfizm
$\mathbb Z\to R$ musi wysłać $[1,0]$ na $1_R$, a stąd
$[n,0]$ na $n1_R$. Z poprzedniej części ma zatem
jedyną możliwą postać. Sprawdzenie mnożenia daje
\[
 (a1_R-b1_R)(c1_R-d1_R)
 =(ac+bd)1_R-(ad+bc)1_R,
\]
co jest obrazem produktu klas. Wzór zachowuje więc
również mnożenie i jedynkę. Jest homomorfizmem pierścieni.
\end{proof}

\begin{przyklad}[Rachunek na klasach formalnych różnic]
\label{aa:prz:z-pairs}
Klasy $[7,3]$ i $[4,0]$ są równe, ponieważ
$7+0=3+4$. Klasa $[2,9]$ odpowiada liczbie $-7$;
ściślej, $[2,9]=[0,7]$, bo $2+7=9+0$.
Mnożąc klasy $[2,5]$ i $[7,3]$, otrzymujemy
$[2,5][7,3]=[29,41]=[0,12]$, czyli $(-3)\cdot4=-12$.
Każda klasa ma postać $[n,0]$ lub $[0,n]$ dla pewnego
$n\in\mathbb N_0$: gdy $a\ge b$, piszemy $a=b+n$,
a gdy $b\ge a$, piszemy $b=a+n$ i stosujemy definicję
relacji $\sim$.
\end{przyklad}

\begin{uwaga}[Dlaczego to nie jest zwykły pierścień ilorazowy?]
\label{aa:uw:z-vs-quotient-ring}
Konstrukcja zaczyna się od półpierścienia $\mathbb N_0$
i relacji równoważności na parach liczb naturalnych;
nie dzielimy istniejącego pierścienia przez ideał.
W języku algebry ogólnej jest to \emph{dopełnienie grupowe}
przemiennego monoidu $(\mathbb N_0,+)$, a poprzednie
twierdzenie wyraża jego własność uniwersalną.
W późniejszej części skryptu skonstruujemy $\mathbb Q$
jako ciało ułamków pierścienia $\mathbb Z$.
Konstrukcję ciała $\mathbb R$ --- w szczególności uzasadnienie
jego zupełności i porządku --- pozostawiamy skryptowi analizy.
Tutaj, przyjmując ciało $\mathbb R$ za dane, skonstruujemy
$\mathbb C$ na dwa algebraiczne sposoby:
jako $\mathbb R[x]/(x^2+1)$ oraz na zbiorze $\mathbb R^2$,
a następnie udowodnimy ich izomorfizm.
\end{uwaga}

\begin{definicja}[Jednostki, dzielniki zera i elementy nilpotentne]
\label{aa:def:pierscien-jednostki}
Niech $R$ będzie pierścieniem z jedynką.
Element $u\in R$ jest \emph{jednostką}, jeśli ma odwrotność
obustronną $v$, tj. $uv=vu=1_R$. Zbiór jednostek oznaczamy $R^\times$.
W niezerowym pierścieniu przemiennym element $a\ne0_R$ jest
\emph{dzielnikiem zera}, jeśli istnieje $b\ne0_R$ z $ab=0_R$.
Element $a$ jest \emph{nilpotentny}, jeśli $a^n=0_R$ dla pewnego
$n\ge1$, a \emph{idempotentny}, jeśli $a^2=a$.
\end{definicja}
\begin{stwierdzenie}[Grupa jednostek]
\label{aa:st:unit-group}
Dla dowolnego pierścienia $R$ jego jednostki tworzą grupę
$R^\times$ względem mnożenia. W pierścieniu przemiennym
niezerowym jednostka nie jest dzielnikiem zera.
\end{stwierdzenie}
\begin{proof}
Dla jednostek $u,v$ zachodzą równości
$(uv)(v^{-1}u^{-1})=(v^{-1}u^{-1})(uv)=1_R$.
Wynika z nich, że $uv$ ma odwrotność obustronną,
zatem jest jednostką. Element $1_R$ należy do $R^\times$,
a wraz z każdą jednostką jej odwrotność także należy do $R^\times$.
Łączność mnożenia jest dziedziczona z pierścienia.
Jeśli $ua=0_R$, mnożenie przez $u^{-1}$ daje
$a=u^{-1}(ua)=0_R$, dlatego jednostka nie jest dzielnikiem zera.
\end{proof}

\begin{definicja}[Dziedzina całkowitości i ciało]
\label{aa:def:dziedzina-cialo}
Niezerowy pierścień przemienny z jedynką, w którym
$ab=0_R$ pociąga $a=0_R$ lub $b=0_R$, nazywamy
\emph{dziedziną całkowitości}. \emph{Ciałem} nazywamy
niezerowy pierścień przemienny z jedynką, w którym każdy
niezerowy element jest jednostką.
\end{definicja}
\begin{lemat}[Skracanie i skończone dziedziny]
\label{aa:lem:skracanie-pierscien}
Niech $R$ będzie dziedziną całkowitości, a $a,b,c\in R$.
Jeśli $a\ne0_R$ i $ab=ac$, to $b=c$.
Każda skończona dziedzina całkowitości jest ciałem.
\end{lemat}
\begin{proof}
Z $ab=ac$ mamy $a(b-c)=0_R$. Brak dzielników zera
oraz $a\ne0_R$ dają $b-c=0_R$. Teraz niech $R$ będzie
skończona, a $a\ne0_R$. Odwzorowanie $m_a:R\to R$,
$m_a(x)=ax$, jest różnowartościowe: jeśli $ax=ay$, to $x=y$
na mocy pierwszej części. Każde różnowartościowe odwzorowanie
skończonego zbioru w siebie jest surjektywne. Istnieje więc
$b\in R$ z $ab=1_R$. Ponieważ $R$ jest przemienny,
$b$ jest odwrotnością $a$. Wszystkie niezerowe elementy są jednostkami.
\end{proof}

\begin{przyklad}[Rozpoznawanie struktur]
\label{aa:prz:pierscien-struktury}
Pierścienie $\mathbb Z$, $\mathbb Q[x]$ i $\mathbb R[x,y]$
są dziedzinami całkowitości. W $\mathbb Z$ jednostkami są
$1,-1$. W pierścieniu $K[x]$ nad ciałem $K$ jednostkami
są dokładnie niezerowe wielomiany stałe: dla niezerowych
wielomianów $f,g$ zachodzi $\deg(fg)=\deg f+\deg g$,
zatem iloczyn równy $1$ wymusza stopnie zero.
Pierścień $M_2(\mathbb R)$ jest nieprzemienny:
$E_{12}E_{21}=E_{11}$, podczas gdy $E_{21}E_{12}=E_{22}$.
Ma także niezerowe macierze, których iloczyn jest macierzą zerową,
na przykład $E_{12}^2=0$. Nie jest więc dziedziną całkowitości.
W $\mathbb Z/6\mathbb Z$ zachodzi $[2][3]=[0]$, lecz
żaden czynnik nie jest zerem. W $\mathbb Z/8\mathbb Z$
klasa $[2]$ jest nilpotentna, bo $[2]^3=[8]=[0]$,
chociaż $[2]\ne[0]$ i $[2]^2=[4]\ne[0]$.
W pierścieniu $\mathbb R\times\mathbb R$ element $(1,0)$
jest idempotentny i jest dzielnikiem zera; jego dopełnienie
$(0,1)$ też jest idempotentne, a ich suma jest jedynką.
\end{przyklad}

\begin{definicja}[Charakterystyka pierścienia]
\label{aa:def:char-ring}
Niech $R$ będzie niezerowym pierścieniem z jedynką.
\emph{Charakterystyką} $R$ jest najmniejsza liczba całkowita
$n>0$ spełniająca $n1_R=0_R$, o ile istnieje;
w przeciwnym razie charakterystyka wynosi $0$.
\end{definicja}
\begin{stwierdzenie}[Charakterystyka dziedziny]
\label{aa:st:char-domain}
Niech $R$ będzie dziedziną całkowitości. Jeśli
$\operatorname{char}R>0$, to $\operatorname{char}R$ jest liczbą pierwszą.
Dla każdej liczby całkowitej
$n\ge2$ pierścień $\mathbb Z/n\mathbb Z$ ma charakterystykę $n$.
\end{stwierdzenie}
\begin{proof}
Jeśli charakterystyka dziedziny wynosi $n=ab$, gdzie
$1<a,b<n$, to $0_R=n1_R=(a1_R)(b1_R)$.
Ponieważ $R$ jest niezerowy, $n\ne1$. Minimalność $n$ zapewnia
$a1_R,b1_R\ne0_R$, co przeczy brakowi dzielników zera.
Zatem $n$ jest liczbą pierwszą.
W $\mathbb Z/n\mathbb Z$ równość $m[1]=[0]$ zachodzi
dokładnie wtedy, gdy $n$ dzieli $m$; najmniejsze dodatnie $m$
to $n$.
\end{proof}

\section{Podpierścienie, homomorfizmy i jądra}
\begin{definicja}[Podpierścień i homomorfizm]
\label{aa:def:ring-hom}
Niech $R$ i $T$ będą pierścieniami z jedynką.
\emph{Podpierścieniem z jedynką} pierścienia $R$ jest podzbiór
$S\subseteq R$ zawierający $1_R$ i zamknięty na dodawanie,
przeciwne elementy i mnożenie. \emph{Homomorfizmem pierścieni}
$\varphi:R\to T$ nazywamy odwzorowanie zachowujące dodawanie,
mnożenie i jedynkę: $\varphi(a+b)=\varphi(a)+\varphi(b)$,
$\varphi(ab)=\varphi(a)\varphi(b)$ i $\varphi(1_R)=1_T$.
Izomorfizm jest bijektywnym homomorfizmem pierścieni.
\end{definicja}
\begin{lemat}[Własności homomorfizmu]
\label{aa:lem:ring-hom-properties}
Niech $R,T$ będą pierścieniami z jedynką,
a $\varphi:R\to T$ homomorfizmem pierścieni.
Wtedy $\varphi(0_R)=0_T$ oraz
$\varphi(-a)=-\varphi(a)$ dla każdego $a\in R$.
Jeżeli $u\in R^\times$, to $\varphi(u)\in T^\times$.
Obraz homomorfizmu jest podpierścieniem $T$ z jedynką.
\end{lemat}
\begin{proof}
Z $\varphi(0_R)=\varphi(0_R+0_R)=2\varphi(0_R)$
po odjęciu $\varphi(0_R)$ otrzymujemy $\varphi(0_R)=0_T$.
Z $\varphi(a)+\varphi(-a)=\varphi(0_R)=0_T$
wynika $\varphi(-a)=-\varphi(a)$. Dla jednostki $u$
i odwrotności $v$ obrazy spełniają
$\varphi(u)\varphi(v)=\varphi(v)\varphi(u)=1_T$.
Obraz jest zamknięty na działania, przeciwny element,
a dzięki zachowaniu jedynki zawiera $1_T$.
\end{proof}

\begin{przyklad}[Podstawienie w wielomianie]
\label{aa:prz:poly-eval}
Niech $K$ będzie ciałem, a $a\in K$. Wzór
$\operatorname{ev}_a:K[x]\to K$, $f\mapsto f(a)$,
określa homomorfizm pierścieni, ponieważ podstawienie
zachowuje sumy, iloczyny i jedynkę. Jego jądro składa się
z wielomianów mających pierwiastek $a$. Dla każdego $k\ge1$ wielomian $x^k-a^k$ jest podzielny
przez $x-a$, ponieważ
$x^k-a^k=(x-a)\sum_{j=0}^{k-1}x^{k-1-j}a^j$.
Stąd dla każdego $f\in K[x]$ różnica $f(x)-f(a)$ jest
wielokrotnością $x-a$. Zatem
$\ker\operatorname{ev}_a=(x-a)$, a obraz jest całym $K$,
ponieważ zawiera wszystkie wielomiany stałe.
Ten przykład będzie za chwilę przypadkiem ogólnego twierdzenia
o izomorfizmie pierścieni.
\end{przyklad}

\section{Ideały jako warunek poprawności konstrukcji ilorazu}
\begin{definicja}[Ideały lewostronne, prawostronne i dwustronne]
\label{aa:def:ideal}
Niech $R$ będzie pierścieniem. Podgrupa addytywna
$I\le(R,+)$ jest \emph{ideałem lewostronnym}, jeśli
$ra\in I$ dla wszystkich $r\in R$, $a\in I$, oraz
\emph{ideałem prawostronnym}, jeśli $ar\in I$ dla tych par.
Jest \emph{ideałem dwustronnym}, jeśli ma obie własności.
W pierścieniu przemiennym te pojęcia są identyczne i mówimy
po prostu \emph{ideał}. Ideał jest \emph{właściwy}, jeśli $I\ne R$.
\end{definicja}
\begin{lemat}[Jądro homomorfizmu jest ideałem]
\label{aa:lem:kernel-ideal}
Niech $R,T$ będą pierścieniami z jedynką, a $\varphi:R\to T$
ich homomorfizmem. Wtedy zbiór
$\ker\varphi=\{a\in R:\varphi(a)=0_T\}$
jest ideałem dwustronnym $R$.
\end{lemat}
\begin{proof}
Ponieważ $\varphi$ zachowuje dodawanie i przeciwne elementy,
$\ker\varphi$ jest podgrupą addytywną. Jeśli
$a\in\ker\varphi$ i $r\in R$, to
$\varphi(ra)=\varphi(r)\varphi(a)=\varphi(r)0_T=0_T$.
Analogicznie $\varphi(ar)=0_T$. Obie własności ideału zachodzą.
\end{proof}

\begin{lemat}[Jedynka w ideale]
\label{aa:lem:ideal-unit}
Niech $R$ będzie pierścieniem z jedynką, a $I\subseteq R$
jego ideałem dwustronnym. Wówczas $I=R$
wtedy i tylko wtedy, gdy $1_R\in I$. Wystarczy też,
aby $I$ zawierał dowolną jednostkę $u\in R^\times$.
\end{lemat}
\begin{proof}
Jeśli $1_R\in I$, to dla każdego $r\in R$ mamy
$r=r1_R\in I$, zatem $I=R$. Jeśli $u\in I$ jest jednostką,
to $1_R=u^{-1}u\in I$. Jeżeli $I=R$, to $1_R\in I$, co dowodzi również implikacji odwrotnej.
\end{proof}

\begin{definicja}[Ideał generowany przez zbiór]
\label{aa:def:generated-ideal}
Niech $R$ będzie pierścieniem przemiennym z jedynką,
a $a_1,\ldots,a_s\in R$ dla pewnego $s\ge1$.
Przez $(a_1,\ldots,a_s)$ oznaczamy zbiór wszystkich sum $r_1a_1+\cdots+r_sa_s$,
$r_i\in R$. Jest to najmniejszy ideał zawierający
$a_1,\ldots,a_s$. Ideał $(a)$ nazywamy \emph{głównym}.
Dla dowolnego zbioru $S\subseteq R$ ideał $(S)$ składa się
ze skończonych sum $\sum_i r_i s_i$, gdzie $r_i\in R$, $s_i\in S$;
dla $S=\varnothing$ jest to ideał zerowy.
\end{definicja}
\begin{proof}
Przedstawiony zbiór zawiera zero, jest zamknięty na sumy,
przeciwne elementy i mnożenie przez dowolny element $R$,
zatem jest ideałem. Zawiera generatory, ponieważ
$a_i=1_Ra_i$. Każdy ideał zawierający te generatory musi
zawierać wszystkie ich wielokrotności i skończone sumy,
więc musi zawierać cały opisany zbiór.
\end{proof}

\begin{definicja}[Suma, przecięcie i iloczyn ideałów]
\label{aa:def:ideal-ops}
Dla ideałów $I,J$ przemiennego pierścienia $R$ definiujemy
$I+J=\{a+b:a\in I,b\in J\}$, $I\cap J$ w zwykłym
sensie teoriomnogościowym oraz \emph{iloczyn ideałów}
$I J$ jako ideał składający się z \emph{wszystkich skończonych
sum} elementów $ab$ z $a\in I$, $b\in J$.
W szczególności nie należy zastępować iloczynu ideałów
przez zbiór samych pojedynczych iloczynów.
\end{definicja}
\begin{stwierdzenie}[Relacje między działaniami na ideałach]
\label{aa:st:ideals-ops}
Niech $R$ będzie pierścieniem przemiennym z jedynką,
a $I,J\subseteq R$ jego ideałami. Wtedy
$IJ\subseteq I\cap J\subseteq I+J$. Ponadto $I+J$
jest najmniejszym ideałem zawierającym $I\cup J$.
\end{stwierdzenie}
\begin{proof}
Każdy składnik $ab$ z $a\in I$, $b\in J$ leży jednocześnie
w $I$ (bo $I$ jest ideałem) i w $J$. Suma takich składników
również leży w obu ideałach. Jeśli $x\in I\cap J$,
to $x=x+0_R\in I+J$. Opis minimalności $I+J$ wynika
z tego, że każdy ideał zawierający $I$ i $J$ jest zamknięty
na ich sumy.
\end{proof}

\begin{twierdzenie}[Pierścień ilorazowy]
\label{aa:tw:ring-quotient}
Niech $I$ będzie ideałem dwustronnym pierścienia $R$.
Na grupie abelowej $R/I$ istnieje pierścień z jedynką,
którego mnożenie jest zadane wzorem
$(a+I)(b+I)=ab+I$. Odwzorowanie
$\pi:R\to R/I$, $a\mapsto a+I$, jest surjektywnym
homomorfizmem pierścieni o jądrze $I$.
\end{twierdzenie}
\begin{proof}
Trzeba przede wszystkim sprawdzić niezależność od reprezentantów.
Jeśli $a'=a+i$ i $b'=b+j$, gdzie $i,j\in I$, to
$a'b'-ab=aj+ib+ij\in I$; pierwsze dwa składniki należą
do $I$ dzięki własnościom ideału dwustronnego,
a trzeci dzięki dowolnej z nich. Zatem $a'b'+I=ab+I$.
Łączność obu działań, przemienność dodawania i oba prawa
rozdzielności przechodzą z $R$ do ilorazu.
Elementem neutralnym dodawania jest $I$,
a jedynką mnożenia $1_R+I$. Projekcja zachowuje oba działania
i jedynkę, zaś $a+I=I$ wtedy i tylko wtedy, gdy $a\in I$.
\end{proof}

\begin{twierdzenie}[Pierwsze twierdzenie o izomorfizmie pierścieni]
\label{aa:tw:ring-isomorphism}
Niech $R,T$ będą pierścieniami z jedynką i niech
$\varphi:R\to T$ będzie homomorfizmem. Wtedy indukowane odwzorowanie
$\overline\varphi:R/\ker\varphi\to\operatorname{im}\varphi$,
$a+\ker\varphi\mapsto\varphi(a)$, jest izomorfizmem pierścieni.
\end{twierdzenie}
\begin{proof}
Jeżeli $a+\ker\varphi=b+\ker\varphi$, to
$a-b\in\ker\varphi$, więc $\varphi(a)=\varphi(b)$;
odwzorowanie jest dobrze określone. Zachowuje sumy,
iloczyny i jedynkę, ponieważ $\varphi$ je zachowuje.
Jest surjektywne z definicji obrazu. Jeśli
$\overline\varphi(a+\ker\varphi)=0_T$, to
$a\in\ker\varphi$, czyli $a+\ker\varphi$ jest zerem
ilorazu. Zatem jest też różnowartościowe.
\end{proof}

\begin{twierdzenie}[Odpowiedniość ideałów przy ilorazie]
\label{aa:tw:ring-ideal-correspondence}
Niech $R$ będzie pierścieniem z jedynką, $I\subseteq R$
jego ideałem dwustronnym i niech
$\pi:R\to R/I$, $\pi(a)=a+I$, będzie projekcją kanoniczną.
Wówczas przyporządkowanie $J\mapsto J/I$ jest bijekcją
między ideałami dwustronnymi $J$ pierścienia $R$
zawierającymi $I$ a ideałami dwustronnymi pierścienia $R/I$.
Odwrotną bijekcją jest $L\mapsto\pi^{-1}(L)$.
Konstrukcja zachowuje zawieranie.
\end{twierdzenie}
\begin{proof}
Przeciwobraz ideału dwustronnego przez homomorfizm pierścieni
jest ideałem dwustronnym: przeciwobraz jest podgrupą addytywną,
a mnożenie elementu przeciwobrazu przez dowolny element
$R$ przenosi się pod $\pi$ na mnożenie elementu $L$
przez element $R/I$. Ponadto $I=\ker\pi\subseteq\pi^{-1}(L)$.
Dla ideału $J$ z $I\subseteq J$ jego obraz $\pi(J)=J/I$
jest ideałem ilorazu, gdyż $\pi$ jest surjektywne.
Jeżeli $a\in\pi^{-1}(\pi(J))$, to $\pi(a)=\pi(j)$
dla pewnego $j\in J$, a więc $a-j\in I\subseteq J$;
stąd $a\in J$. Z drugiej strony $\pi(\pi^{-1}(L))=L$
z surjektywności $\pi$. Otrzymaliśmy obie odwrotne bijekcje.
Jeśli $J_1\subseteq J_2$, to każda klasa elementu $J_1$
należy również do $J_2/I$, zatem $J_1/I\subseteq J_2/I$.
Analogicznie z $L_1\subseteq L_2$ wynika
$\pi^{-1}(L_1)\subseteq\pi^{-1}(L_2)$.
\end{proof}

\begin{przyklad}[Wszystkie ideały w $\mathbb Z/12\mathbb Z$]
\label{aa:prz:ideals-z12}
Wszystkie ideały pierścienia $\mathbb Z/12\mathbb Z$ można przedstawić
przez przeciwobrazy w $\mathbb Z$. Uzasadnijmy najpierw, że każdy ideał $J\subseteq\mathbb Z$ ma postać
$d\mathbb Z$. Jeśli $J=\{0\}$, bierzemy $d=0$. W przeciwnym
razie w $J$ istnieje dodatnia liczba całkowita, gdyż dla
$z\in J\setminus\{0\}$ również $-z\in J$,
a jedna z liczb $z,-z$ jest dodatnia. Wybieramy najmniejszą
taką liczbę $d$. Dla dowolnego $a\in J$ dzielenie z resztą daje
$a=qd+r$ z $0\le r<d$. Ponieważ $r=a-qd\in J$,
minimalność $d$ wymusza $r=0$. Zatem $J=d\mathbb Z$.
Teraz zawieranie $12\mathbb Z\subseteq d\mathbb Z$ oznacza
$d\mid12$. Warunek ten wyklucza $d=0$.
Dla dodatniego dzielnika $d$ obraz ideału $d\mathbb Z$ ma
$12/d$ elementów, reprezentowanych przez klasy
$[0],[d],\ldots,[(12/d-1)d]$. Zatem mamy dokładnie sześć ideałów:
$\{0\}$, $([6])$, $([4])$, $([3])$, $([2])$ i cały pierścień.
Na przykład $([4])=\{[0],[4],[8]\}$, a
$([6])=\{[0],[6]\}$. Ich indeksy i liczności wynikają
z dzielenia $12$ przez odpowiednie dzielniki.
\end{przyklad}

\section{Ideały pierwsze i maksymalne}
\begin{definicja}[Ideał pierwszy i maksymalny]
\label{aa:def:prime-max-ideal}
Niech $R$ będzie pierścieniem przemiennym z jedynką.
Ideał właściwy $P\subsetneq R$ jest \emph{pierwszy}, gdy
dla wszystkich $a,b\in R$ warunek $ab\in P$ pociąga
$a\in P$ lub $b\in P$.
Ideał właściwy $M\subsetneq R$ jest \emph{maksymalny}, gdy
nie istnieje ideał $J$ spełniający
$M\subsetneq J\subsetneq R$.
\end{definicja}
\begin{twierdzenie}[Kryteria przez pierścień ilorazowy]
\label{aa:tw:prime-max-quotient}
Niech $R$ będzie pierścieniem przemiennym z jedynką,
a $I\subsetneq R$ jego ideałem właściwym.
Wówczas $I$ jest pierwszy wtedy i tylko wtedy, gdy
$R/I$ jest dziedziną całkowitości. Ponadto $I$ jest
maksymalny wtedy i tylko wtedy, gdy $R/I$ jest ciałem.
W szczególności każdy ideał maksymalny jest pierwszy.
\end{twierdzenie}
\begin{proof}
W ilorazie $R/I$ mamy $(a+I)(b+I)=I$ dokładnie wtedy,
gdy $ab\in I$. Klasa $a+I$ jest zerem dokładnie wtedy,
gdy $a\in I$. Ponieważ $I$ jest właściwy, iloraz nie jest
pierścieniem zerowym. Warunek braku dzielników zera jest
zatem równoważny definicji ideału pierwszego.

Jeżeli $I$ jest maksymalny, to z twierdzenia o odpowiedniości
ideałów wynika, że w $R/I$ jedynymi ideałami są $\{0\}$
i cały pierścień. Dla dowolnego $a+I\ne0$ ideał
$(a+I)$ nie jest zerowy, więc jest całym $R/I$.
Istnieje zatem klasa $b+I$ z $(a+I)(b+I)=1_R+I$,
czyli każda klasa niezerowa jest odwracalna. Iloraz jest ciałem.
Odwrotnie, w ciele każdy ideał niezerowy zawiera element
odwracalny, a więc jest całym ciałem. Twierdzenie
o odpowiedniości daje maksymalność $I$.
Wreszcie każde ciało jest dziedziną całkowitości:
z $ab=0$ i $a\ne0$ wynika $b=a^{-1}(ab)=0$.
\end{proof}

\begin{przyklad}[Ideały w $\mathbb Z$ i wielomiany]
\label{aa:prz:prime-max-examples}
Dla liczby pierwszej $p$ ideał $p\mathbb Z$ jest
maksymalny, ponieważ $\mathbb Z/p\mathbb Z$ jest ciałem.
Ideał $(0)$ w $\mathbb Z$ jest pierwszy, ale nie maksymalny:
iloraz $\mathbb Z/(0)$ to $\mathbb Z$, który nie jest ciałem.
W $K[x]$ ideał $(x-a)$ jest maksymalny dla $a\in K$,
ponieważ przez podstawienie w $a$ mamy izomorfizm
$K[x]/(x-a)\cong K$. Natomiast w $K[x,y]$ ideał $(x)$ jest pierwszy,
ale nie maksymalny. Rzeczywiście, podstawienie $x=0$
określa surjektywny homomorfizm $K[x,y]\to K[y]$ o jądrze
$(x)$: każdy wielomian ma postać $f(0,y)+xg(x,y)$.
Z pierwszego twierdzenia o izomorfizmie wynika
$K[x,y]/(x)\cong K[y]$, a ten ostatni pierścień
jest dziedziną całkowitości, lecz nie jest ciałem.
\end{przyklad}

\section{Chińskie twierdzenie o resztach dla pierścieni}
\begin{definicja}[Ideały względnie pierwsze]
\label{aa:def:comaximal}
Dwa ideały $I,J$ przemiennego pierścienia $R$ nazywamy
\emph{względnie pierwszymi}, jeśli $I+J=R$.
Jest to równoważne istnieniu $u\in I$ i $v\in J$ z $u+v=1_R$.
\end{definicja}
\begin{lemat}[Przecięcie i iloczyn ideałów względnie pierwszych]
\label{aa:lem:comaximal-intersection}
Niech $R$ będzie pierścieniem przemiennym z jedynką,
a $I,J\subseteq R$ ideałami. Jeśli $I+J=R$, to $I\cap J=IJ$.
\end{lemat}
\begin{proof}
Inkluzja $IJ\subseteq I\cap J$ została już udowodniona.
Niech $u\in I$, $v\in J$ oraz $u+v=1_R$.
Jeśli $z\in I\cap J$, to $z=zu+zv$.
Pierwszy składnik leży w $IJ$, ponieważ $z\in J$, $u\in I$;
drugi też leży w $IJ$, ponieważ $z\in I$, $v\in J$.
Zatem $z\in IJ$ i otrzymujemy równość.
\end{proof}
\begin{twierdzenie}[Chińskie twierdzenie o resztach]
\label{aa:tw:crt-ring}
Niech $R$ będzie pierścieniem przemiennym z jedynką,
a $I_1,\ldots,I_s$ parami względnie pierwszymi ideałami
($s\ge2$). Wówczas homomorfizm
\[
\Phi:R\longrightarrow \prod_{i=1}^{s}R/I_i,
\qquad r\longmapsto(r+I_1,\ldots,r+I_s)
\]
jest surjektywny, ma jądro $I_1\cap\cdots\cap I_s$
i indukuje izomorfizm pierścieni
\[
R/(I_1\cdots I_s)\ \cong\ \prod_{i=1}^{s}R/I_i.
\]
\end{twierdzenie}
\begin{proof}
Jądro $\Phi$ to elementy należące do wszystkich ideałów,
czyli ich przecięcie. Udowodnimy surjektywność jawnie.
Dla $i\ne j$ wybierzmy $u_{ij}\in I_i$ i $v_{ij}\in I_j$
takie, że $u_{ij}+v_{ij}=1_R$. Dla ustalonego $i$
połóżmy $e_i=\prod_{j\ne i}v_{ij}$. Mamy wtedy
$e_i\equiv1_R\pmod{I_i}$, ponieważ każdy czynnik
$v_{ij}=1_R-u_{ij}$ ma tę własność, oraz
$e_i\equiv0_R\pmod{I_j}$ dla każdego $j\ne i$, bo w
iloczynie występuje czynnik $v_{ij}\in I_j$.
Dla dowolnie wybranych klas $r_i+I_i$ element
$r=\sum_i r_i e_i$ ma obraz $\Phi(r)=(r_i+I_i)_i$.
To dowodzi surjektywności.

Pozostaje pokazać równość przecięcia z iloczynem dla wielu
ideałów. Jeśli $I+J=R$ i $I+K=R$, wybieramy
$u+v=1_R$ i $u'+w=1_R$ z $u,u'\in I$, $v\in J$, $w\in K$.
Mnożenie daje $1_R=uu'+uw+vu'+vw$, przy czym
pierwsze trzy składniki należą do $I$, a ostatni do $JK$.
Zatem $I+JK=R$. Stosując ten fakt wielokrotnie, otrzymujemy,
że $I_s$ jest względnie pierwszy z $I_1\cdots I_{s-1}$.
Indukcja i poprzedni lemat pokazują
$I_1\cap\cdots\cap I_s=I_1\cdots I_s$.
Pierwsze twierdzenie o izomorfizmie pierścieni kończy dowód.
\end{proof}

\begin{przyklad}[Rachunek modulo $12$ i idempotenty]
\label{aa:prz:crt12}
Ideały $3\mathbb Z$ i $4\mathbb Z$ są względnie pierwsze,
ponieważ $4-3=1$. Chińskie twierdzenie o resztach daje
$\mathbb Z/12\mathbb Z\cong\mathbb Z/3\mathbb Z\times\mathbb Z/4\mathbb Z$.
Wprost: klasa $[4]_{12}$ przechodzi na parę $([1]_3,[0]_4)$,
a klasa $[9]_{12}$ na $([0]_3,[1]_4)$.
Obie są idempotentne, bo $4^2\equiv4\pmod{12}$ i
$9^2\equiv9\pmod{12}$, a ich iloczyn jest zerowy.
Dla układu $x\equiv2\pmod3$, $x\equiv3\pmod4$ dostajemy
$x\equiv 2\cdot4+3\cdot9=35\equiv11\pmod{12}$.
Wynik można sprawdzić: $11-2=9$ jest podzielne przez $3$,
a $11-3=8$ jest podzielne przez $4$.
Natomiast dla ideałów $2\mathbb Z$ i $4\mathbb Z$
warunek względnej pierwszości nie zachodzi; układ
$x\equiv0\pmod2$, $x\equiv1\pmod4$ nie ma rozwiązania.
\end{przyklad}

\begin{uwaga}[Od pierścieni do geometrii i modułów]
\label{aa:uw:ring-next}
Ideały pierścieni wielomianów opisują układy równań
algebraicznych, a ilorazy pierścieni pozwalają zachować
informację nie tylko o zbiorach zer, lecz również o krotnościach
i elementach nilpotentnych. Moduły z poprzedniego rozdziału
są strukturami addytywnymi, na których działa pierścień.
W kolejnym rozdziale sprawdzimy, kiedy podzielność w pierścieniu
jest na tyle dobrze zorganizowana, aby każdy ideał był główny.
\end{uwaga}

\setlength{\parskip}{4pt}
